\section{Đối chiếu yêu cầu bổ sung của Assignment 2}
Mười task Assignment 1 và khung ba Sprint của phần chính được giữ nguyên. Assignment 2 bổ sung năm task quy trình, mỗi task tối đa hai điểm, triển khai qua hai Sprint theo đề; không thay Product Backlog chức năng hoặc tự cộng hai thang điểm thành điểm môn học.

\begin{table}[htbp]\centering\small
\begin{tabularx}{\textwidth}{|c|X|X|}\hline
\textbf{Task A2} & \textbf{Nhóm yêu cầu} & \textbf{Tạo phẩm dùng để đối chiếu}\\\hline
1 & Bối cảnh, SRS, ưu tiên và use-case & Đặc tả yêu cầu, quy tắc nghiệp vụ, AC và UC-01.\\\hline
2 & Mô hình hóa UML & Use-case, activity, sequence, class; nguồn và hình render.\\\hline
3 & Kiến trúc và triển khai & Kiến trúc phân lớp, quyết định ADR, mô hình môi trường và ranh giới dịch vụ.\\\hline
4 & Quản lý repo, tài liệu và CI & Nhánh/PR, Definition of Done, các bước install/lint/test/build.\\\hline
5 & Delivery và nghiệm thu & Bộ test, đóng gói, triển khai, migration, smoke, rollback, release và demo.\\\hline
\end{tabularx}\caption{Các nhóm tạo phẩm đối chiếu với Assignment 2}
\end{table}

\subsection{Kiến trúc và quyết định cần ghi nhận}
Mô tả kiến trúc cần nối được Next.js/PWA, các module NestJS, PostgreSQL/migration và payment service. ADR ghi bối cảnh, lựa chọn, phương án khác, ưu/nhược và tác động. Những lựa chọn đã có trong code như monorepo, Prisma, gateway tách riêng, idempotency/snapshot cần được ghi nhận; ngày viết ADR không tự chứng minh quyết định đã được review từ trước.

Deployment diagram phải thể hiện thiết bị, frontend/CDN, API, PostgreSQL persistent storage, registry, HTTPS/cổng, secret, health/log và ranh giới local/CI/production. Tên thành phần trong mô hình được ánh xạ sang thư mục frontend/backend của repo để truy vết tới mã nguồn.

\subsection{CI, CD và điều kiện bàn giao}
CI bắt buộc tự chạy khi mở PR: install, lint, unit test và build; nhánh chính yêu cầu checks và không merge khi thất bại. Cần lưu link PR cùng log/ảnh một lần pass và một lần bắt lỗi. Commit viết ``test passed'' không tự chứng minh đã có pipeline.

CD cần quality gate, image gắn SHA/version, registry nếu nền tảng yêu cầu, deploy bằng secret CI, migration an toàn và health/smoke sau deploy. Thất bại phải dừng, có hướng dẫn rollback app và cách khôi phục dữ liệu; không dùng reset hoặc seed đè production. Bàn giao gồm HTTPS frontend, API/health, database không public trực tiếp, release/tag, release notes, README local/test/build/deploy/rollback, demo và Review/Retrospective.

\subsection{Truy vết và giới hạn bằng chứng}
Ma trận \texttt{docs/traceability-matrix.md} nối FR/NFR với use-case, thiết kế, code/API và test. Ma trận giữ liên kết hai chiều giữa yêu cầu, mô hình, code và test theo từng phiên bản tài liệu. Nguồn chuẩn của SRS và quy tắc nằm ở docs, bản report dùng cùng mã yêu cầu.

\section{Đặc tả chi tiết và tiêu chí chấp nhận}
\begingroup\footnotesize\setlength{\tabcolsep}{3pt}
\begin{longtable}{|>{\raggedright\arraybackslash}p{0.07\textwidth}|>{\raggedright\arraybackslash}p{0.2\textwidth}|>{\raggedright\arraybackslash}p{0.07\textwidth}|>{\raggedright\arraybackslash}p{0.58\textwidth}|}
\caption{Yêu cầu chức năng và tiêu chí chấp nhận}\\\hline
\textbf{Mã} & \textbf{Yêu cầu} & \textbf{Ưu tiên} & \textbf{Tiêu chí chấp nhận có thể kiểm chứng} \\ \hline
\endfirsthead
\hline
\textbf{Mã} & \textbf{Yêu cầu} & \textbf{Ưu tiên} & \textbf{Tiêu chí chấp nhận có thể kiểm chứng} \\ \hline
\endhead
FR-01 & Xem menu& Must & AC-01: GET /products trả danh sách có id, tên, giá, ảnh; UI có loading, empty và error; lỗi tải không làm mất điều hướng \\ \hline
FR-02 & Xem chi tiết và tùy chọn& Must & AC-02: GET /products/:id trả chi tiết/options; chọn S/M/L, topping cập nhật giá; id không tồn tại có trạng thái không tìm thấy \\ \hline
FR-03 & Giỏ hàng& Must & AC-03: thêm/gộp đúng tùy chọn, đổi số lượng, xóa; tổng bằng tổng dòng; reload còn giỏ; giỏ rỗng không gửi đơn \\ \hline
FR-04 & Đăng ký, đăng nhập và JWT& Must & AC-04: đăng ký email duy nhất; sai mật khẩu không cấp token; API bảo vệ thiếu/sai JWT trả 401; response không lộ passwordHash \\ \hline
FR-05 & Tạo đơn và tính giá ở backend& Must & AC-05: input hợp lệ trả Order PENDING; tiền từ quy tắc backend; sai size/topping/qty bị từ chối không đổi database; cùng key/nội dung không tạo thêm đơn \\ \hline
FR-06 & Thanh toán Ví/Thẻ mô phỏng& Must & AC-06: khi chỉ bật mock, thành công ghi Payment/PAID; lỗi có kết quả PAYMENT\_FAILED; mock bị chặn khi VNPay bật để không bỏ qua cổng \\ \hline
FR-07 & Xác nhận, theo dõi và lịch sử& Must & AC-07: chủ đơn xem đúng mã/trạng thái/món/tiền; GET /orders/me không trả đơn người khác; refresh/realtime không nhân đôi dữ liệu \\ \hline
FR-08 & Barista quản lý trạng thái& Must & AC-08: chỉ BARISTA gọi staff API; chuyển trạng thái đúng máy trạng thái đơn hàng; không pha đơn thiếu tiền hoặc có REQUESTED; trạng thái cũ bị từ chối khi thao tác đồng thời \\ \hline
FR-09 & Giữ kho đồng thời và hủy an toàn độ đúng nghiệp vụ & Should & AC-09: hai khách tranh phần cuối chỉ một đơn giữ thành công; kho không âm; định lượng đúng size/topping; hủy hợp lệ hoàn snapshot một lần dù Recipe đã đổi; đơn hết hạn giữ món được đóng và thông báo đặt lại; đơn đã pha không được tự hủy \\ \hline
FR-10 & Voucher và điểm thưởng & Should & AC-10: mã hết hạn/inactive bị từ chối; giảm và điểm theo chính sách giảm giá và tích điểm; callback lặp không cộng lại điểm \\ \hline
FR-11 & VNPay URL và QR & Should & AC-11: backend tạo URL ký đúng số tiền còn phải trả; QR mở cùng URL và có nút mở trang; URL hết hạn/đơn không payable không hiển thị QR khả dụng \\ \hline
FR-12 & Kết quả cổng an toàn & Should & AC-12: sai chữ ký/merchant/amount không đổi PAID; Return chỉ chuyển hướng; IPN/QueryDr hợp lệ cập nhật đúng một lần trong transaction \\ \hline
FR-13 & Thanh toán lại, kết quả muộn và thu trùng & Should & AC-13: cùng key trả cùng lần; mã còn hạn tái dùng; sau hết hạn QR và đơn còn hạn giữ món có thể cấp lần mới, không xóa lần cũ; hai lần thành công chỉ một lần áp dụng cho đơn và khoản dư có hồ sơ xử lý \\ \hline
FR-14 & Khách yêu cầu đổi món & Should & AC-14: liên kết hỗ trợ nhỏ trên đơn PAID mở menu/giỏ riêng; chọn size/topping bình thường; báo giá từ backend; cùng key không tạo trùng; REQUESTED có thể rút bằng version \\ \hline
FR-15 & Quầy duyệt đổi và thu chênh lệch & Should & AC-15: giá đã xác nhận giữ theo snapshot; duyệt chuyển credit/đổi giữ kho nguyên tử; tiền thu mới = total - credit; thấp hơn giá trị cũ bị từ chối; hai người cùng duyệt chỉ một thành công \\ \hline
FR-16 & Theo dõi hoàn khoản thu không áp dụng & Should & AC-16: thu dư mở một hồ sơ theo attempt; BARISTA nhận xử lý, nhập chứng từ rồi hoàn tất theo version; chứng từ trùng trả lỗi; khách không tự xác nhận hoàn \\ \hline
FR-17 & Hàng đợi hỗ trợ tại quầy & Should & AC-17: queue lọc trạng thái, tìm mã/email, phân trang 20 dòng; xem món cũ/mới, tiền và ghi chú; thiếu quyền 403; page/status sai 400 \\ \hline
FR-18 & Đóng gói và chạy local& Must & AC-18: clone sạch, điền biến mẫu, một lệnh Compose chạy web/API/DB; migration+seed trên DB sạch tạo được dữ liệu; health endpoint phản ánh tình trạng dịch vụ \\ \hline
FR-19 & CI/CD và bàn giao& Must & AC-19: PR tự install/lint/test/build; failed check không merge; release tạo image có version/SHA, deploy, migrate an toàn, smoke test; có cách rollback và URL HTTPS frontend/API/health \\ \hline
FR-20 & Giao hàng & Won't & AC-20: phạm vi bàn giao ghi rõ chưa có UI tài xế/điều phối/tracking; không quảng bá trạng thái hiện tại là vòng đời giao hàng đầy đủ \\ \hline
\end{longtable}
\endgroup
\subsection{Yêu cầu phi chức năng và phép kiểm chứng}
Các ngưỡng là mục tiêu nghiệm thu, chưa phải số liệu đo đã đạt.
\begingroup\footnotesize\setlength{\tabcolsep}{3pt}
\begin{longtable}{|>{\raggedright\arraybackslash}p{0.07\textwidth}|>{\raggedright\arraybackslash}p{0.16\textwidth}|>{\raggedright\arraybackslash}p{0.07\textwidth}|>{\raggedright\arraybackslash}p{0.62\textwidth}|}
\caption{NFR và phép kiểm chứng}\\\hline
\textbf{Mã} & \textbf{Yêu cầu} & \textbf{Ưu tiên} & \textbf{Tiêu chí chấp nhận / phép đo} \\ \hline
\endfirsthead
\hline
\textbf{Mã} & \textbf{Yêu cầu} & \textbf{Ưu tiên} & \textbf{Tiêu chí chấp nhận / phép đo} \\ \hline
\endhead
NFR-01 & Bảo mật xác thực/phân quyền & Must & Mật khẩu băm; khách không đọc đơn người khác/gọi staff; 401/403/404 đúng hợp đồng; không có secret trong Git/log/ảnh \\ \hline
NFR-02 & Toàn vẹn tiền và kho & Must & Chạy test cạnh tranh/idempotency/rollback: không kho âm, không nhân đôi settlement/điểm, không chuyển credit vượt tổng; mọi lỗi giữa transaction không để dữ liệu dở dang \\ \hline
NFR-03 & Hiệu năng API nội bộ & Should & Với 50 sản phẩm, 20 người dùng đồng thời, sau warm-up: p95 GET sản phẩm và GET đơn \(\leq\)500ms trên môi trường ghi rõ CPU/RAM/DB; tách độ trễ gateway khỏi số liệu; ghi tỷ lệ lỗi và lệnh đo \\ \hline
NFR-04 & Mobile-first và khả năng truy cập & Must & Ở 390px và 1440px không tràn ngang; nhãn form, focus bàn phím, thông báo lỗi và trạng thái chờ rõ; nút chính tối thiểu 44px; xác nhận tiền không chỉ dựa màu sắc \\ \hline
NFR-05 & Tính khả dụng và phục hồi & Must & Có loading/empty/error, retry có kiểm soát; mất Return không mất kết quả đã ghi; không treo UI vô hạn; health check production dùng cho smoke \\ \hline
NFR-06 & Tái tạo và triển khai & Must & Install từ lockfile; build frontend/backend; migration+seed trên database mới; Compose đủ 3 dịch vụ; HTTPS và DB production không public trực tiếp \\ \hline
NFR-07 & Quan sát và hỗ trợ & Must & Log có thời gian và tham chiếu lỗi, không lộ token/secret; phân biệt Order/attempt/refund; có hướng dẫn xử lý lỗi thanh toán; health không trả cấu hình nhạy cảm \\ \hline
NFR-08 & Bảo trì và truy vết & Must & Mỗi FR/NFR liên kết use-case, thiết kế, code/test; sơ đồ có nguồn+render; quyết định có ADR; không có đường dẫn/file minh chứng giả \\ \hline
NFR-09 & Quality gate & Must & CI tự chạy trên PR, branch protection yêu cầu checks; lưu link/log một lần thành công và một lần phát hiện lỗi; test không dùng database production \\ \hline
NFR-10 & Sao lưu, migration và khôi phục & Must & Backup trước migration; thử khôi phục DB riêng; script deploy không reset hoặc seed đè DB đang dùng; runbook rollback nêu phần app và phần dữ liệu \\ \hline
NFR-11 & Phiên đăng nhập nhất quán & Should & Chính sách một tài khoản/origin hoặc phiên riêng/tab được chốt; đổi tài khoản/logout không để response cũ ghi đè user/token mới; có test hai tab và reload \\ \hline
NFR-12 & Cài đặt PWA & Could & Có manifest, icon và kiểm tra cài đặt nếu nhóm chọn triển khai PWA; không coi hỗ trợ mobile là bằng chứng đã có offline/PWA \\ \hline
\end{longtable}
\endgroup

\section{Đặc tả chi tiết UC-01}
\begingroup\footnotesize\setlength{\tabcolsep}{3pt}
\begin{longtable}{|>{\raggedright\arraybackslash}p{0.24\textwidth}|>{\raggedright\arraybackslash}p{0.67\textwidth}|}
\caption{Thông tin UC-01}\\\hline
\textbf{Thuộc tính} & \textbf{Nội dung} \\ \hline
\endfirsthead
\hline
\textbf{Thuộc tính} & \textbf{Nội dung} \\ \hline
\endhead
Actor chính & Khách hàng \\ \hline
Actor hỗ trợ & Dịch vụ thanh toán; barista nhận đơn đủ tiền và hỗ trợ khi có bất thường \\ \hline
Trigger & Khách bấm Đặt hàng từ giỏ có món, sau đó chủ động bắt đầu thanh toán \\ \hline
Mục tiêu & Có một đơn đúng giá/tùy chọn, giữ nguyên liệu hợp lệ và xác định kết quả thu tiền; khách nhận được mã/trạng thái đơn \\ \hline
Tiền điều kiện & Menu và cấu hình giá tồn tại; khách đăng nhập trước khi gửi tạo đơn; PostgreSQL sẵn sàng; gateway mock hoặc VNPay được cấu hình theo chế độ \\ \hline
Hậu điều kiện thành công & Order PAID, reservation còn giữ; đúng một khoản thanh toán áp dụng cho đơn; điểm trên tiền mới được ghi cùng transaction; đơn có thể xuất hiện tại quầy \\ \hline
Hậu điều kiện tối thiểu & Không kho âm/đơn nửa chừng/thu lặp do callback; giao dịch chưa rõ giữ lịch sử để đối soát; frontend không tự coi timeout là thất bại chắc chắn \\ \hline
Liên kết yêu cầu & FR-01--08, FR-09--13; NFR-01/02/04/05/07/11 \\ \hline
Quy tắc & Giá do backend chốt; giữ/hoàn theo snapshot; thanh toán xác nhận một lần; chỉ chuyển trạng thái hợp lệ \\ \hline
\end{longtable}
\endgroup
\begingroup\footnotesize\setlength{\tabcolsep}{3pt}
\begin{longtable}{|>{\raggedright\arraybackslash}p{0.07\textwidth}|>{\raggedright\arraybackslash}p{0.35\textwidth}|>{\raggedright\arraybackslash}p{0.49\textwidth}|}
\caption{Luồng chính UC-01}\\\hline
\textbf{Bước} & \textbf{Khách / Next.js UI} & \textbf{NestJS / database / dịch vụ thanh toán} \\ \hline
\endfirsthead
\hline
\textbf{Bước} & \textbf{Khách / Next.js UI} & \textbf{NestJS / database / dịch vụ thanh toán} \\ \hline
\endhead
M1 & Mở menu, xem chi tiết & GET /products và /products/:id trả sản phẩm và tùy chọn \\ \hline
M2 & Chọn size/topping/số lượng, thêm giỏ & Chưa có Order, chưa khóa/giữ kho \\ \hline
M3 & Mở giỏ, xác nhận món, voucher nếu có & UI tính tiền tham khảo; backend mới là nguồn giá chính thức \\ \hline
M4 & Đăng nhập nếu cần, gửi tạo đơn với key cho lần đặt & JWT và DTO được kiểm tra; key được nhận trước khi giữ nguyên liệu \\ \hline
M5 & Chờ kết quả tạo đơn & Transaction tính lại tiền/voucher, khóa nguyên liệu có thứ tự, kiểm tra hàng, trừ lượng giữ, tạo Order PENDING và snapshot; lỗi thì rollback \\ \hline
M6 & Nhận mã đơn và bấm Thanh toán & Backend xác thực chủ đơn, trạng thái và snapshot; khóa Order; tạo/tái dùng attempt theo key; số tiền từ total - creditApplied \\ \hline
M7 & Quét QR mở URL hoặc bấm mở trang VNPay & VNPay Sandbox nhận URL ký; mock thực hiện theo chế độ kiểm thử riêng khi VNPay tắt \\ \hline
M8 & Thực hiện xác nhận trên cổng & IPN đến backend; backend xác minh chữ ký, merchant/ref/amount/status trước khi nhận kết quả \\ \hline
M9 & Có thể quay về trang đơn qua Return & Return chuyển về giao diện; trình duyệt đọc trạng thái backend. IPN/QueryDr mới xác nhận Payment, PAID và điểm trong transaction \\ \hline
M10 & Xem xác nhận, lịch sử và tiến độ & UI đọc trạng thái từ backend, cập nhật qua HTTP/realtime; barista bắt đầu pha nếu đơn đủ điều kiện \\ \hline
\end{longtable}
\endgroup
\begingroup\footnotesize\setlength{\tabcolsep}{3pt}
\begin{longtable}{|>{\raggedright\arraybackslash}p{0.18\textwidth}|>{\raggedright\arraybackslash}p{0.44\textwidth}|>{\raggedright\arraybackslash}p{0.29\textwidth}|}
\caption{Luồng thay thế UC-01}\\\hline
\textbf{Mã / điểm rẽ} & \textbf{Luồng} & \textbf{Kết quả} \\ \hline
\endfirsthead
\hline
\textbf{Mã / điểm rẽ} & \textbf{Luồng} & \textbf{Kết quả} \\ \hline
\endhead
A1 tại M4 & Chưa đăng nhập \(\rightarrow\) mở đăng nhập với next phù hợp \(\rightarrow\) quay về giỏ & Không tạo đơn trước xác thực \\ \hline
A2 tại M4--M5 & Mất response \(\rightarrow\) gửi lại cùng key/nội dung & Trả lại cùng Order; không trừ kho thêm \\ \hline
A3 tại M6 & Cùng key tạo thanh toán hoặc key khác khi mã còn hiệu lực & Trả/tái dùng attempt còn hiệu lực, không tạo nhiều mã hoạt động do bấm lặp \\ \hline
A4 tại M7--M9 & Xác nhận giao dịch thất bại & Lưu kết quả; đơn có thể PAYMENT\_FAILED, nguyên liệu chưa tự hoàn chỉ vì thanh toán lỗi; khách retry hoặc hủy nếu đủ điều kiện \\ \hline
A5 tại M7 & Mã hết hạn, chưa có khoản thu đã xác minh & Backend thử QueryDr, có thể cấp attempt mới sau hết hạn; giữ lần cũ để nhận kết quả muộn, cảnh báo khách nguy cơ thu trùng \\ \hline
A6 tại M9 & Return đến trước IPN hoặc khách đóng trình duyệt & Đơn chưa tự PAID từ Return; mở lịch sử/kiểm tra kết quả, IPN/QueryDr có thể cập nhật sau \\ \hline
A7 sau M5 & Khách muốn hủy đơn chưa trả & Kiểm tra đơn chưa thanh toán và còn được hủy; nếu hợp lệ hủy và hoàn snapshot cùng transaction \\ \hline
A8 sau M10 & Khách cần đổi món chưa pha & Chuyển UC-09; không hủy đơn cũ khi khách mới mở menu đổi \\ \hline
\end{longtable}
\endgroup
\begingroup\footnotesize\setlength{\tabcolsep}{3pt}
\begin{longtable}{|>{\raggedright\arraybackslash}p{0.07\textwidth}|>{\raggedright\arraybackslash}p{0.33\textwidth}|>{\raggedright\arraybackslash}p{0.51\textwidth}|}
\caption{Ngoại lệ UC-01}\\\hline
\textbf{Mã} & \textbf{Điều kiện} & \textbf{Xử lý và dữ liệu cần giữ} \\ \hline
\endfirsthead
\hline
\textbf{Mã} & \textbf{Điều kiện} & \textbf{Xử lý và dữ liệu cần giữ} \\ \hline
\endhead
E1 & Size/topping/qty/product/voucher không hợp lệ & Trả lỗi có thể hiểu; không tạo đơn/giữ kho dở dang \\ \hline
E2 & Hai khách tranh lượng cuối, người sau không đủ nguyên liệu & Một transaction thành công, transaction thiếu hàng rollback; thông báo thiếu món; không stock âm \\ \hline
E3 & Cùng key tạo đơn nhưng khác nội dung & 409, yêu cầu kiểm tra đơn trước; không tự dùng key đó cho đơn khác \\ \hline
E4 & Sai JWT hoặc không sở hữu đơn & Từ chối theo guard/quyền; không trả dữ liệu người khác \\ \hline
E5 & Callback sai chữ ký/ref/merchant/amount & Không cập nhật PAID/Payment/điểm; phản hồi theo giao thức cổng \\ \hline
E6 & QueryDr timeout/không có kết quả chắc chắn & Giữ trạng thái chưa rõ; không suy ra ngân hàng đã trả tiền; UI có đường kiểm tra/lịch sử/hỗ trợ \\ \hline
E7 & Callback lặp, đồng thời hoặc đến sai thứ tự & Unique/khóa Order/idempotency; không hạ thanh toán đã xác nhận thành thất bại do callback cũ \\ \hline
E8 & Lỗi ghi Payment hoặc điểm & Rollback transaction; callback có thể được gửi lại và xử lý an toàn \\ \hline
E9 & Hai attempt đều thu thành công & Một khoản áp dụng cho đơn, khoản còn lại có REVIEW/paidAt để quầy xử lý hoàn khoản thu dư \\ \hline
E10 & Tiền về sau khi đơn đã đóng/không còn nguyên liệu giữ hợp lệ & Không tự pha/trừ kho lại; giữ khoản thu và hồ sơ hỗ trợ \\ \hline
E11 & Ngrok/Return cấu hình sai hoặc cảnh báo trình duyệt & Return có thể chưa về UI; kết quả phải kiểm tra từ backend; không thanh toán lại chỉ vì chưa chuyển trang \\ \hline
\end{longtable}
\endgroup
\begingroup\footnotesize\setlength{\tabcolsep}{3pt}
\begin{longtable}{|>{\raggedright\arraybackslash}p{0.13\textwidth}|>{\raggedright\arraybackslash}p{0.33\textwidth}|>{\raggedright\arraybackslash}p{0.32\textwidth}|>{\raggedright\arraybackslash}p{0.13\textwidth}|}
\caption{Kịch bản nghiệm thu UC-01}\\\hline
\textbf{ID} & \textbf{Tiền đề / bước kiểm tra} & \textbf{Mong đợi} & \textbf{Yêu cầu} \\ \hline
\endfirsthead
\hline
\textbf{ID} & \textbf{Tiền đề / bước kiểm tra} & \textbf{Mong đợi} & \textbf{Yêu cầu} \\ \hline
\endhead
UC01-T01 & Menu \(\rightarrow\) tùy chọn \(\rightarrow\) giỏ \(\rightarrow\) đăng nhập \(\rightarrow\) tạo \(\rightarrow\) mock thành công \(\rightarrow\) xác nhận & Một đơn, giá đúng, PAID, có lịch sử & FR-01--07 \\ \hline
UC01-T02 & Hai yêu cầu đồng thời mua phần cuối & Một thành công, một thiếu hàng; kho không âm & FR-09, NFR-02 \\ \hline
UC01-T03 & Gửi lại tạo đơn và IPN nhiều lần & Một Order cho key, một Payment áp dụng, điểm một lần & FR-05/12 \\ \hline
UC01-T04 & Mock thất bại rồi retry; hủy đơn chưa trả nếu đủ điều kiện & Trạng thái/giữ kho/hoàn snapshot theo điều kiện hủy đơn chưa thanh toán & FR-06/09 \\ \hline
UC01-T05 & VNPay Return hợp lệ nhưng chưa IPN & Không tự PAID; giao diện chờ/kiểm tra & FR-12 \\ \hline
UC01-T06 & VNPay IPN sai chữ ký hoặc sai tiền & Không ghi nhận thanh toán & FR-12, NFR-01 \\ \hline
UC01-T07 & Mã cũ hết hạn, query lỗi, mã mới được cấp, cả hai IPN thành công & Một lần phục vụ, khoản dư có hồ sơ & FR-13/16 \\ \hline
UC01-T08 & Đổi 45.000đ sang 75.000đ, duyệt rồi thu thêm & UI nổi bật 30.000đ; backend chỉ thu 30.000đ; không trừ kho lần hai khi trả thêm & FR-14/15 \\ \hline
\end{longtable}
\endgroup

